\PhanI
%\loai{Dao động tắt dần}
\begin{ex} %[3P1-21.1Y][Loai1] 
(MH - 2018) Dao động cơ tắt dần
\choice
{có biên độ tăng dần theo thời gian}
{\True có biên độ giảm dần theo thời gian}
{luôn có hại}
{luôn có lợi}
\loigiai{
Dao động tắt dần có biên độ giảm dần theo thời gian}
\end{ex} \haidong{20}
\begin{ex}%[3P1-19.2B][Loai1]
Một vật dao động tắt dần có các đại lượng giảm liên tục theo thời gian là
\choice
{li độ và tốc độ}
{biên độ và gia tốc}
{biên độ và tốc độ}
{\True biên độ và năng lượng}
\loigiai{
Một vật dao động tắt dần thì biên độ và năng lượng của vật giảm dần theo thời gian}
\end{ex} \haidong{20}
\begin{ex}%[3P1-19.2B][Loai1]
Nguyên nhân gây ra dao động tắt dần của con lắc đơn trong không khí là do
\choice
{dây treo có khối lượng đáng kể}
{lực căng dây treo}
{trọng lực tác dụng lên vật}
{\True lực cản môi trường}
\loigiai{
Nguyên nhân gây ra dao động tắt dần của con lắc đơn trong không khí là do lực cản môi trường}
\end{ex} \haidong{20}
\begin{ex}%[3P1-19.2B][Loai1]
Khi nói về dao động cơ tắt dần của một vật, phát biểu nào sau đây đúng?
\choice
{Động năng của vật biến thiên theo hàm bậc nhất của thời gian}
{Cơ năng của vật không thay đổi theo thời gian}
{Lực cản của môi trường tác dụng lên vật càng nhỏ thì dao động tắt dần càng nhanh}
{\True Biên độ dao động của vật giảm dần theo thời gian}
\loigiai{
\begin{itemize}
\item A. Sai vì động năng của vật biến thiên theo hàm sin hoặc cos theo thời gian.
\item B. Sai vì trong dao động tắt dần biên độ của vật sẽ thay đổi nên cơ năng của vật cũng bị thay đổi theo.
\item C. Sai vì khi lực cản của môi trường tác dụng lên vật càng lớn thì dao động tắt dần càng nhanh và ngược lại.
\item D. Đúng vì trong dao động tắt dần biên độ của vật sẽ giảm dần theo thời gian.
\end{itemize}}
\end{ex} \haidong{20}

\begin{ex}%[3P1-21.1K][Loai1]
Phát biểu nào sau đây là \textbf{sai}? Đối với dao động tắt dần thì
\choice
{cơ năng giảm dần theo thời gian}
{\True tần số giảm dần theo thời gian}
{ma sát và lực cản càng lớn thì dao động tắt dần càng nhanh}
{biên độ dao động giảm dần theo thời gian}
\loigiai{
Đối với dao động tắt dần thì tần số giảm dần theo thời gian là SAI}
\end{ex} \haidong{20}
\begin{ex}
(QG - 2021) Khi một con lắc lò xo đang dao động tắt dần do tác dụng của lực ma sát thì cơ năng của con lắc chuyển hóa dần dần thành
\choice
{điện năng}
{thế năng}
{\True nhiệt năng}
{hóa năng}
\end{ex} \haidong{20}
\begin{ex}%[3P1-21.1B][Loai1]
Một vật dao động cưỡng bức dưới tác dụng của một ngoại lực biến thiên điều hòa với tần số $f$. Chu kì dao động của vật là
\choice
{$f$}
{$\dfrac{2}{f}$}
{\True $\dfrac{1}{f}$}
{$2f$}
\loigiai{
Tần số dao động của vật bằng với tần số của ngoại lực tác dụng lên vật và bằng $f$. Nên chu kì dao động của vật là: $T=\dfrac{1}{f}$}
\end{ex} \haidong{20}
\begin{itemize}
\item[]
\end{itemize}
\begin{ex}
Trong dao động tắt dần một phần cơ năng đã biến đổi thành
\choice
{điện năng}
{\True nhiệt năng}
{hoá năng}
{quang năng}
\end{ex} \haidong{20}
%\loai{Dao động cưỡng bức}
\begin{ex} 
(QG - 2022) Một hệ dao động cưỡng bức, phát biểu nào sau đây \textbf{sai}?
\choice
{Dao động cưỡng bức có biên độ không đổi}
{Dao động cưỡng bức có biên độ phụ thuộc vào biên độ của lực cưỡng bức}
{\True Dao động cưỡng bức có tần số luôn bằng tần số dao đông riêng của hệ}
{Dao động cưỡng bức có tần số bằng tần số của lực cưỡng bức}
\end{ex} \haidong{20}
\begin{ex}%[3P1-21.1B][Loai1]
Khi nói về dao động cơ cưỡng bức, phát biểu nào sau đây là \textbf{sai}?
\choice
{Biên độ của dao động cưỡng bức phụ thuộc vào tần số của lực cưỡng bức}
{Dao động cưỡng bức có tần số luôn bằng tần số của lực cưỡng bức}
{Biên độ của dao động cưỡng bức phụ thuộc vào biên độ của lực cưỡng bức}
{\True Dao động cưỡng bức có tần số luôn bằng tần số riêng của hệ dao động}
\loigiai{
Dao động cưỡng bức có tần số luôn bằng tần số của ngoại lực cưỡng bức}
\end{ex} \haidong{20}
\begin{ex}%[3P1-21.1B][Loai1]
Biên độ dao động cưỡng bức \textbf{không} phụ thuộc vào
\choice
{hệ số lực cản tác dụng lên vật}
{biên độ ngoại lực tuần hoàn tác dụng lên vật}
{\True pha ban đầu của ngoại lực tuần hoàn tác dụng lên vật}
{tần số ngoại lực tuần hoàn tác dụng lên vật}
\loigiai{
Biên độ của dao động cưỡng bức phụ thuộc vào:
\begin{itemize}
\item Tần số ngoại lực tuần hoàn tác dụng lên vật và tần số riêng của hệ.
\item Biên độ ngoại lực tuần hoàn tác dụng lên vật.
\item Lực ma sát (lực cản) của môi trường.
\end{itemize}}
\end{ex} \haidong{20}

\begin{ex}%[3P1-21.1K][Loai1]
Khi nói về dao động cưỡng bức ở giai đoạn ổn định, phát biểu nào sau đây là \textbf{sai}?
\choice
{Biên độ dao động phụ thuộc vào tần số của ngoại lực cưỡng bức}
{Tần số dao động bằng tần số của ngoại lực}
{Biên độ dao động phụ thuộc vào ma sát}
{\True Tần số ngoại lực tăng thì biên độ dao động tăng}
\loigiai{
Khi nói về dao động cưỡng bức ở giai đoạn ổn định tần số ngoại lực tăng thì biên độ dao động tăng là SAI}
\end{ex} \haidong{20}
\begin{ex}%[3P1-21.1B][Loai1]
Phát biểu nào sau đây là \textbf{sai}?
\choice
{Sự cộng hưởng thể hiện rõ nét nhất khi lực ma sát của môi trường ngoài là nhỏ}
{\True Biên độ cộng hưởng không phụ thuộc vào ma sát}
{Biên độ dao động cưỡng bức phụ thuộc vào mối quan hệ giữa tần số của lực cưỡng bức và tần số dao động riêng của hệ}
{Dao động cưỡng bức là dao động dưới tác dụng của ngoại lực biến đổi tuần hoàn}
\loigiai{
Biên độ cộng hưởng phụ thuộc vào ma sát, cụ thể lực ma sát càng nhỏ thì biên độ khi xảy ra cộng hưởng càng lớn}
\end{ex} \haidong{20}
%\loai{Hiện tượng cộng hưởng}
\begin{ex} %[3P1-19.1Y][Loai1] 
(QG - 2018)Một con lắc lò xo có tần số dao động riêng $f_0$. Khi tác dụng vào nó một ngoại lực cưỡng bức tuần hoàn có tần số f thì xảy ra hiện tượng cộng hưởng. Hệ thức nào sau đây đúng?
\choice
{\True $f = f_0$}
{$f = 4f_0$}
{$f = 0,5f_0$}
{$f = 2f_0$}
\loigiai{
Hiện tượng cộng hưởng xảy ra khi: $f = f_0$.
}
\end{ex} \haidong{20}
\begin{ex}%[3P1-21.1K][Loai1]
Hiện tượng cộng hưởng thể hiện rõ nét khi
\choice
{tần số lực cưỡng bức nhỏ}
{biên độ lực cưỡng bức nhỏ}
{\True lực cản môi trường nhỏ}
{tần số lực cưỡng bức lớn}
\loigiai{
Hiện tượng cộng hưởng thể hiện rõ nét khi lực cản môi trường nhỏ}
\end{ex} \haidong{20}
\begin{ex}%[3P1-21.1B][Loai1]
Trong hiện tượng cộng hưởng thì
\choice
{biên độ ngoại lực cưỡng bức đạt cực đại}
{\True biên độ dao động cưỡng bức đạt cực đại}
{tần số dao động cưỡng bức đạt cực đại}
{tần số dao động riêng đạt giá trị cực đại}
\loigiai{
Trong hiện tượng cộng hưởng thì biên độ của dao động cưỡng bức đạt cực đại}
\end{ex} \haidong{20}
\begin{ex}%[3P1-21.1B][Loai1]
Hiện tượng cộng hưởng cơ học xảy ra khi
\choice
{tần số của lực cưỡng bức bằng tần số của dao động cưỡng bức}
{\True tần số dao động cưỡng bức bằng tần số dao động riêng của hệ}
{tần số của lực cưỡng bức lớn hơn tần số riêng của hệ}
{tần số của lực cưỡng bức bé hơn tần số riêng của hệ}
\loigiai{
Hiện tượng cộng hưởng xảy ra khi tần số của dao động cưỡng bức bằng với tần số dao động riêng của hệ}
\end{ex} \haidong{20}


\begin{ex}%[3P1-21.1B][Loai1]
\nguon{MH - 2019} Thực hiện thí nghiệm về dao động cưỡng bức như hình sau.\\
\centerline{\begin{tikzpicture}[draw=Mapcolor]
\tkzDefPoints{0/0/a, 5/0/b, 5/4/c, 0/4/d, 1.5/3/m}
\tkzDefPointBy[homothety=center m ratio 0.5](c)\tkzGetPoint{1}
\tkzDefPointBy[homothety=center m ratio 0.6](c)\tkzGetPoint{2}
\tkzDefPointBy[homothety=center m ratio 0.7](c)\tkzGetPoint{3}
\tkzDefPointBy[homothety=center m ratio 0.8](c)\tkzGetPoint{4}
\tkzDefShiftPoint[m](-90:2cm){m'}
\tkzDefShiftPoint[1](-90:1.8cm){1'}
\tkzDefShiftPoint[2](-90:0.9cm){2'}
\tkzDefShiftPoint[3](-90:1.5cm){3'}
\tkzDefShiftPoint[4](-90:0.6cm){4'}
\tkzDrawSegments[line width=2pt, blue](a,b b,c a,d)
\tkzDrawSegments[line width=1pt, blue](c,m d,m m,m' 1,1' 2,2' 3,3' 4,4')
\tkzDrawSegments[dashed,red](c,d)
\node [line width=2pt,draw,fill=Mapcolor!60, thick, shape=rectangle, minimum width=0.8cm, minimum height=0.6cm, anchor=center] at (0,0.2){\tiny $  $};
\node [line width=2pt,draw,fill=Mapcolor,pattern=north east lines, thick, shape=rectangle, minimum width=0.8cm, minimum height=0.6cm, anchor=center] at (0,0.2){\tiny $  $};
\node [line width=2pt,draw,fill=Mapcolor!60, thick, shape=rectangle, minimum width=0.8cm, minimum height=0.6cm, anchor=center] at (5,0.2){\tiny $  $};
\node [line width=2pt,draw,fill=Mapcolor,pattern=north east lines, thick, shape=rectangle, minimum width=0.8cm, minimum height=0.6cm, anchor=center] at (5,0.2){\tiny $  $};
\shade[ball color=Mapcolor] (m') circle (1.3ex)node[below=5pt]{$M$};
\shade[ball color=Mapcolor] (1') circle (1ex)node[below=5pt]{$(1)$};
\shade[ball color=Mapcolor] (2') circle (1ex)node[below=5pt]{$(2)$};
\shade[ball color=Mapcolor] (3') circle (1ex)node[below=5pt]{$(3)$};
\shade[ball color=Mapcolor] (4') circle (1ex)node[below=5pt]{$(4)$};
\end{tikzpicture}}
Năm con lắc đơn: $(1)$, $(2)$, $(3)$, $(4)$ và M (con lắc điều khiển) được treo trên một sợi dây. Ban đầu hệ đang đứng yên ở vị trí cân bằng. Kích thích M dao động nhỏ trong mặt phẳng vuông góc với mặt phẳng hình vẽ thì các con lắc còn lại dao động theo. Không kể M, con lắc dao động mạnh nhất là
\choice
{con lắc $(2)$}
{\True con lắc $(1)$}
{con lắc $(3)$}
{con lắc $(4)$}
\loigiai{
Nhận thấy con lắc $(1)$ có chiều dài gần bằng con lắc M (con lắc điều khiển) tức là tần số dao động gần bằng nhau nên con lắc $(1)$ sẽ dao động mạnh nhất. Đây là ứng dụng của hiện tượng cộng hưởng cơ}
\end{ex} \haidong{20}
\PhanII
\setcounter{ex}{0}
\begin{ex}
Cho các nhận định sau về dao động tắt dần.
\choiceTF[t]
{Dao động tắt dần có động năng giảm dần còn thế năng biến thiên điều hòa}
{\True Dao động tắt dần là dao động có biên độ giảm dần theo thời gian}
{\True Lực ma sát càng lớn thì dao động tắt càng nhanh}
{\True Trong dao động tắt dần, cơ năng giảm dần theo thời gian}
\loigiai{
\begin{itemize}
\item Sai.
\item Đúng.
\item Đúng.
\item Đúng.
\end{itemize}
}
\end{ex} \haidong{20}


\begin{ex}
Cho các nhận định sau về một con lắc dao động cưỡng bức.
\choiceTF[t]
{Dao động cưỡng bức có tần số luôn bằng tần số dao động riêng}
{Biên độ dao động cưỡng bức không phụ thuộc vào biên độ ngoại lực}
{Biên độ dao động cưỡng bức luôn phụ thuộc vào pha của ngoại lực}
{\True Dao động cưỡng bức có tần số bằng tần số của ngoại lực}
\loigiai{
\begin{itemchoice}
\itemch
\itemch
\itemch
\itemch
\end{itemchoice}
}
\end{ex} \haidong{20}

\begin{ex}
Một hệ dao động điều hòa với tần số dao động riêng $4 \ Hz$. Tác dụng vào hệ dao động đó một ngoại lực có biểu thức $F=F_0 \cos (8 \pi t+\dfrac{\pi}{3})\ N$.
\choiceTF[t]
{Hệ sẽ dao động cưỡng bức với tần số dao động là $8 \ Hz$}
{\True Hệ sẽ dao động với biên độ cực đại vì đang xảy ra hiện tượng cộng hưởng}
{Hệ sẽ ngừng dao động vì do hiệu tần số của ngoại lực cưỡg bức và tần số dao động riêng bằng 0}
{Hệ sẽ dao động với biên độ giảm dần rất nhanh do ngoại lực tác dụng cản trở dao động}
\loigiai{
\begin{itemchoice}
\itemch
\itemch
\itemch
\itemch
\end{itemchoice}
}
\end{ex} \haidong{20}
\begin{ex}
Cho các nhận định sau về hiện tượng cộng hưởng trong dao động cưỡng bức.
\choiceTF[t]
{\True Đường cong cộng hưởng càng nhọn khi lực cản của môi trường càng nhỏ}
{\True Hiện tượng cộng hưởng là hiện tượng biên độ dao động cưỡng bức tăng đến giá trị cực đại khi tần số ngoại lực bằng tần số dao động riêng của hệ}
{Trong kĩ thuật, hiện tượng cộng hưởng cơ học luôn gây ra nhiều tác hại}
{\True Các công trình như những cây cầu dài hay nhà cao tầng thường được thiết kế để phòng tránh hiện tượng cộng hưởng xảy ra}
\loigiai{
\begin{itemchoice}
\itemch
\itemch
\itemch
\itemch
\end{itemchoice}
}
\end{ex} \haidong{20}
